\section{Lazy versus rich parameterization}
\label{sec:lazy}

Whether $K(t)$ stays at $K(0)$ is a scaling question~\cite{chizat2019lazy,yang2021mup,pehlevan2023infinite}. Write
\begin{align}
  h^{(1)}_\mu &= N^{-a_1} W^{(1)} x_\mu, &
  h^{(\ell+1)}_\mu &= N^{-a_{\ell+1}} W^{(\ell+1)}\ph\!\big(h^{(\ell)}_\mu\big), &
  f_\mu &= \gamma^{-1} N^{-a_L}\, w\cdot\ph\!\big(h^{(L-1)}_\mu\big),
  \label{eq:pb-mlp}
\end{align}
with $W^{(\ell)}_{ij}\sim\mathcal{N}(0,N^{-b_\ell})$ and learning rate
\begin{equation}
  \eta=\eta_0\,\gamma^2 N^{c},
  \qquad
  \gamma=\gamma_0 N^{d},\qquad \gamma_0=\ON{1}.
  \label{eq:lr-scale}
\end{equation}
Backpropagated fields:
\begin{equation}
  g^{(\ell)}_{\mu,i}
  \equiv
  N^{a_{\ell}+\frac12 b_{\ell}}\,
  \pd{f_\mu}{h^{(\ell)}_{\mu,i}},
  \qquad
  G^{(\ell)}_{\mu\nu}
  \equiv
  N^{-1}\, g^{(\ell)}_\mu\cdot g^{(\ell)}_\nu.
  \label{eq:G-kernel}
\end{equation}

\paragraph{Preactivations $\ON{1}$ at initialization.}
For $\ell\ge 2$,
\begin{equation}
  \E\bigl[\lvert h^{(\ell)}_{\mu i}\rvert^2\bigr]
  =N^{-2a_\ell-b_\ell}\,
  \E\bigl[\lvert\ph(h^{(\ell-1)}_\mu)\rvert^2\bigr].
  \label{eq:h-var}
\end{equation}
A finite nonzero limit requires
\begin{equation}
  2a_\ell+b_\ell=1,\quad \ell=2,\ldots,L,
  \qquad
  2a_1+b_1=0.
  \label{eq:init-constraint}
\end{equation}
Under~\eqref{eq:init-constraint} the hidden units of a layer become i.i.d.\ Gaussians and $\Phi^{(\ell)}$ concentrates on
\begin{equation}
  \Phi^{(\ell)}_{\mu\nu}
  =\E_{h\sim\mathcal{N}(0,\Phi^{(\ell-1)})}\bigl[\ph(h_\mu)\ph(h_\nu)\bigr],
  \qquad
  \Phi^{(0)}_{\mu\nu}=D^{-1}\,x_\mu\cdot x_\nu.
  \label{eq:phi-rec-init}
\end{equation}

\paragraph{Predictions move in $\ON{1}$ time.}
From~\eqref{eq:fn-ode}, $\partial_t f=\ON{1}$ requires $\eta K/P=\ON{1}$. Expanding $K$ over layers as in~\eqref{eq:two-ntk} and using $\Phi,G=\ON{1}$,
\begin{equation}
\begin{split}
  \frac{\gamma^2}{N^{c}}\,\partial_\theta f_\mu\cdot\partial_\theta f_\nu
  &=
  N^{-c-2a_L+1}\Phi^{(L-1)}_{\mu\nu}
  +\sum_{\ell=2}^{L-1} N^{-c-2a_\ell+1} G^{(\ell)}_{\mu\nu}\Phi^{(\ell-1)}_{\mu\nu}\\
  &\qquad
  +N^{-c-2a_1} G^{(1)}_{\mu\nu}\Phi^{(0)}_{\mu\nu}.
\end{split}
  \label{eq:K-expand}
\end{equation}
Every exponent vanishes iff
\begin{equation}
  2a_\ell+c=1\quad(\ell\ge 2),
  \qquad
  2a_1+c=0.
  \label{eq:pred-constraint}
\end{equation}

\paragraph{Features move, or they do not.}
Gradient flow on $W^{(1)}$ gives
\begin{equation}
  \partial_t h^{(1)}_\mu
  =N^{-2a_1-c+d-\frac12}\,
  \frac{\eta_0\gamma_0}{P}\sum_{\nu}\Delta_\nu\, g^{(1)}_\nu\,\Phi^{(0)}_{\mu\nu}.
  \label{eq:h1-dot}
\end{equation}
Imposing $\partial_t h^{(1)}=\ON{1}$ together with~\eqref{eq:pred-constraint} forces $d=\tfrac12$. The same exponent appears at every later layer, because
\begin{equation}
  \partial_t h^{(\ell)}_\mu
  =N^{-2a_\ell-c+d-\frac12}\times(\ON{1})
  +N^{-a_\ell}W^{(\ell)}\bigl(\dot\ph(h^{(\ell-1)}_\mu)\odot\partial_t h^{(\ell-1)}_\mu\bigr)
  \label{eq:hl-dot}
\end{equation}
and the second term inherits the scaling of $\partial_t h^{(\ell-1)}$. Any $d<\tfrac12$ sends $\partial_t h\to 0$: features freeze and $K(t)\to K(0)$. The NTK parameterization of Jacot et al.\ is $d=0$.

\begin{proposition}
The family $d=\tfrac12$,~\eqref{eq:init-constraint},~\eqref{eq:pred-constraint} yields infinite-width feature learning. A representative (mean-field / $\mu$P) is
\begin{equation}
  a_1=0,\quad
  a_{\ell\ge 2}=\tfrac12,\quad
  b_\ell=0,\quad
  c=0,\quad
  d=\tfrac12,
  \qquad
  \eta=\eta_0\gamma_0^2 N,
  \qquad
  f=\frac{w\cdot\ph(h^{(L-1)})}{\gamma_0\sqrt{N}}.
  \label{eq:mup}
\end{equation}
Then $f_\mu(0)=\ON{N^{-1/2}}$. The NTK parameterization is the same architecture with $d=0$, or equivalently $\gamma_0\to 0$ after a time rescaling.
\end{proposition}

Chizat's diagnostic: training $\alpha f(\theta/\alpha)$ as $\alpha\to\infty$ linearizes about initialization. Here $\alpha\sim\gamma_0^{-1}N^{1/2-d}$.

\begin{figure}[ht]
\centering
\includegraphics[width=0.95\textwidth]{fig02_lazy_rich.pdf}
\caption{Two-layer ReLU, $N=256$, $\eta=\eta_0\gamma_0^2 N$. Left: relative motion of the hidden feature kernel $\Phi$. Right: training square loss. Only $\gamma_0$ is varied.}
\label{fig:lazy-rich}
\end{figure}

\begin{exercise}
For $f=\gamma_0^{-1}N^{-1/2}v\cdot\ph(Wx/\sqrt{D})$ and $c=0$, compute the $N$-scaling of $\partial_t W$ and $\partial_t v$. Recover $d=\tfrac12$ as the condition that $W$ moves by $\ON{1}$.
\end{exercise}
